%% start of file `template.tex'.
%% Copyright 2006-2015 Xavier Danaux (xdanaux@gmail.com), 2020-2024 moderncv maintainers (github.com/moderncv).
%
% This work may be distributed and/or modified under the
% conditions of the LaTeX Project Public License version 1.3c,
% available at http://www.latex-project.org/lppl/.


\documentclass[11pt,a4paper,sans]{moderncv}        % possible options include font size ('10pt', '11pt' and '12pt'), paper size ('a4paper', 'letterpaper', 'a5paper', 'legalpaper', 'executivepaper' and 'landscape') and font family ('sans' and 'roman')

\moderncvcolor{black}
\moderncvstyle[]{banking}
\nopagenumbers{}                                  % uncomment to suppress automatic page numbering for CVs longer than one page

\usepackage{hyperref}

% adjust the page margins
\usepackage[hmargin=0.5in,vmargin=10pt]{geometry}

% font loading
% for luatex and xetex, do not use inputenc and fontenc
% see https://tex.stackexchange.com/a/496643
\ifxetexorluatex
  \usepackage{fontspec}
  \usepackage{unicode-math}
  \usepackage{FiraSans}
  \defaultfontfeatures{Ligatures=TeX}
  \setmainfont{Fira Sans}
  \setsansfont[ItalicFont={Fira Sans Light Italic},
            BoldFont={Fira Sans Bold},
            BoldItalicFont={Fira Sans Italic}]
            {Fira Sans Light}
  % \setmonofont{Latin Modern Mono}
  \setmonofont{Monaspace Neon}
  \setmathfont{Latin Modern Math}

  % you may also consider Fira Sans Light for a extra modern look
  %\setsansfont[ItalicFont={Fira Sans Light Italic},%
  %           BoldFont={Fira Sans},%
  %           BoldItalicFont={Fira Sans Italic}]%
  %           {Fira Sans Light}%
\else
  \usepackage[utf8]{inputenc}
  \usepackage[T1]{fontenc}
  \usepackage{lmodern}
\fi

% document language
\usepackage[UTF8,scheme=plain]{ctex}  % FIXME: using spanish breaks moderncv
% \setCJKmainfont{SimHei}
% \setCJKsansfont{SimHei}
\setCJKmainfont{Noto Sans CJK SC}
\setCJKsansfont{Noto Sans CJK SC}

\usepackage{setspace}
\linespread{1.5}

% personal data
\name{\textmd{\textsc{吴胜桥}}}{}
\phone[mobile]{+86~159~8583~0350}
\email{Blind-Guess-Senior@outlook.com}
\homepage{https://blind-guess-senior.github.io/}

\social[github]{Blind-Guess-Senior}

% \extrainfo{additional information}                 % optional, remove / comment the line if not wanted
% \quote{Some quote}                                 % optional, remove / comment the line if not wanted

%----------------------------------------------------------------------------------
%            content
%----------------------------------------------------------------------------------
\begin{document}
% \begin{CJK*}{UTF8}{gbsn}                          % to typeset your resume in Chinese using CJK
%-----       resume       ---------------------------------------------------------
\makecvtitle

\section{教育经历}  % for 'contemporary' style use optional argument for displaying an icon, e.g. \section[\faGraduationCap]{Education}
\cventry{2024 -- 至今}{本科 | \textit{计算机科学与技术}}{上海科技大学}{\textit{中国 上海}}{}{}  % arguments 3 to 6 can be left empty
% \cventry{year--year}{Degree}{Institution}{City}{\textit{Grade}}{Description}

\section{专业技能}
\cvitem{语言}{C\# | C/C++ | Javascript/Typescript | Nix | OCaml | Python | LaTeX}

\section{比赛经历}
\cventry{2022.10}{CCF 非专业级软件能力认证提高组（CSP-S） 第二轮}{CCF CSP-S 第二轮 \textmd{\textsc{| 二等奖}}}{C++, Argorithm}{}{}
\cventry{2022.10}{CCF 非专业级软件能力认证提高组（CSP-S） 第一轮}{CCF CSP-S 第一轮 \textmd{\textsc{| 二等奖}}}{C++, Argorithm}{}{}

\section{游戏项目}
\cventry{2026.07}{\textcolor[RGB]{128,128,128}{以“拍键盘”为核心输入的2D节奏游戏}}{KBre \textmd{\textsc{| 团队项目}}}{Unity}{}{
  \vspace{0.2em}
  主程 | 5 人团队
  \vspace{0.1em}
  \begin{itemize}
    \item 玩家以“拍键盘”的方式控制角色移动、攻击。拍击、滑动、连续拍击等输入模式有各自的特殊效果。玩家需要配合游戏节拍输入以提高攻击伤害并降低资源消耗
    \item 设计整体代码框架，设计并实现事件系统、输入模式、场景物体、场景控制、敌人状态机等
  \end{itemize}}
\vspace{0.4em}

\cventry{2025.07}{\textcolor[RGB]{128,128,128}{2D能力锁平台跳跃游戏}}{C-17ES-α \textmd{\textsc{| 课程作业}}}{Unity}{}{
  \vspace{0.2em}
  主程 | 5 人团队
  \vspace{0.1em}
  \begin{itemize}
    \item 玩家初始仅有向右的移动能力，随流程逐步解锁向左、向上、向下方向的移动能力。不同方向的移动具有独特的效果，玩家需要组合运用这些能力，跨越障碍并完成关卡
    \item 设计整体代码框架，设计并实现事件系统、移动能力、场景内交互物体等
  \end{itemize}}
\vspace{0.4em}

\section{项目经历}
\cventry{2026.03 -- 至今}{\textcolor[RGB]{128,128,128}{基于预编译对话脚本的 Unity 对话系统库}}{Artifact Dialoguer \textmd{\textsc{| 个人项目}}}{C\#, .NET, Unity}{}{
  \vspace{0.2em}
  \begin{itemize}
    \item 设计简洁且无歧义的对话脚本语法，并实现 LL(1) 语法解析器将其编译为 Unity ScriptableObject
    \item 实现 Unity 内的运行时对话控制器，读取编译产物并驱动游戏内对话流程
  \end{itemize}
}

\cventry{2026.02 -- 至今}{\textcolor[RGB]{128,128,128}{轻量 Unity 工具库}}{Artifact Unity Utils \textmd{\textsc{| 个人项目}}}{C\#, .NET, Unity}{}{
  \vspace{0.2em}
  \begin{itemize}
    \item 基于 C\# 和 .NET 实现事件总线、依赖注入、有限状态机等通用架构与四叉树等数据结构
  \end{itemize}
}
\vspace{0.4em}

\cventry{2025.09 -- 2026.01}{\textcolor[RGB]{128,128,128}{OATC语言编译器}}{OATC Language Compiler \textmd{\textsc{| 个人项目}}}{OCaml, LLVM, X86lite}{}{
  \vspace{0.2em}
  \begin{itemize}
    \item 使用 OCaml 实现 X86lite 指令集模拟器、汇编器，及 OATC 语言解释器
    \item 实现从 OATC 到 LLVMlite IR 的编译器前端，以及从 LLVMlite IR 到 X86lite 的编译器后端，从而实现完整编译流程
  \end{itemize}
}
\vspace{0.4em}

\cventry{2025.02 -- 2025.06}{\textcolor[RGB]{128,128,128}{面向80x86架构的操作系统}}{PintOS \textmd{\textsc{| 个人项目}}}{C, x86 Assembly}{}{
  \vspace{0.2em}
  \begin{itemize}
    \item 实现高级调度、系统调用、用户态/内核态机制、虚拟内存及文件系统等操作系统核心模块
    \item 最终的 PintOS 包含基础 shell 和文件系统，且支持在用户态运行程序
  \end{itemize}
}
\vspace{0.4em}
% \vspace{0.5em}

% Detailed achievements:
% \begin{itemize}
% \item Achievement 1
% \item Achievement 2 (with sub-achievements)
%   \begin{itemize}
%   \item Sub-achievement (a);
%   \item Sub-achievement (b), with sub-sub-achievements (don't do this!);
%     \begin{itemize}
%     \item Sub-sub-achievement i;
%     \item Sub-sub-achievement ii;
%     \item Sub-sub-achievement iii;
%     \end{itemize}
%   \item Sub-achievement (c);
%   \end{itemize}
% \item Achievement 3
% \item Achievement 4
% \end{itemize}}
% \cventry{year--year}{Job title}{Employer}{City}{}{Description line 1\newline{}Description line 2\newline{}Description line 3}
% \subsection{Miscellaneous}
% \cventry{year--year}{Job title}{Employer}{City}{}{Description}

\clearpage

\end{document}


%% end of file `template.tex'.
